\section{Full Tables}\label{sec:full-tables}

Full tables provide more ways to define table values than simple tables. In a simple
table, the table value is obtained by interpolation. Values of the independent
variables---the state variables---are known from flight-path integration, and the
tabulated data values are known from the input. Linear interpolation is used to
compute the table value, or dependent variable. The only flexibility is in the
selection of the independent variables and their values.

In a full table, the table value is obtained by executing a sequence of math
operations. When TAOS needs a table value, the value is first initialized to zero.
Then the operations supplied in the table input are executed one at a time. Each
operation adds, subtracts, multiplies, or otherwise combines the current table value
with another value. The other value can be a constant, a value interpolated from a
table, a state variable, or a user-defined variable. The process is analogous to using
a calculator with an accumulation register.

The sequence of math operations resembles a simple programming language. The
operations let the user control how a table value is computed. Temporary values can
be stored for later use, and simple \texttt{if-then} and \texttt{goto} statements
provide flow control. The following example illustrates the concept:

\begin{lstlisting}[style=taosmanual]
(1st-stage)
  table thrust units=lb
  start
    sub pres          # subtract atmospheric pressure
    mult 3.219        # multiply by nozzle exit area
    add tvac(time)    # add the vacuum thrust
      time = 0, 0.1, 0.2, 0.5, 1.0, 4.0, 4.1, 4.2, 4.3
      tvac = 0, 1200, 1050, 980, 950, 910, 100, 50, 0
  end
\end{lstlisting}

This table defines thrust magnitude in pounds force as vacuum thrust minus
atmospheric pressure multiplied by nozzle exit area:

\begin{equation}
  \left\lVert\vec{F}_{\mathrm{prop}}\right\rVert
  = \left\lVert\vec{F}_{\mathrm{vac}}\right\rVert - pS_{\mathrm{noz}}.
  \label{eq:full-table-thrust}
\end{equation}

The calculation begins at \texttt{start}, which initializes thrust to zero. The
statement \texttt{sub pres} subtracts atmospheric pressure from zero, leaving the
negative pressure as the current table value. The statement \texttt{mult 3.219}
multiplies that value by a constant representing nozzle exit area. Finally,
\texttt{add tvac(time)} interpolates vacuum thrust as a function of time and adds it
to the current table value. The result is used in the flight-path calculations.

As with any programming language, there are many ways to perform the same
calculation. One could instead multiply pressure by exit area, save the result in a
temporary variable, clear the table value, interpolate vacuum thrust, and subtract the
saved pressure-area product. That method requires more operations and is less
efficient.

Tables are evaluated many times while a flight path is being calculated, so efficient
table definitions are worthwhile. Reducing the number of operations can
significantly reduce processing time.

\subsection{Math Operations}\label{sec:table-math-operations}

Math operations available in TAOS are listed in \cref{tab:math-operations}. A
sequence always begins with \texttt{start} and ends with the \texttt{end} operation.
The last operation in a table must be \texttt{end}.

\begin{table}[htbp]
  \centering
  \caption{Math operations.}
  \label{tab:math-operations}
  \small
  \begin{tabularx}{\linewidth}{@{}>{\ttfamily}lX>{\ttfamily}lX@{}}
    \toprule
    \normalfont Operation & Description & \normalfont Operation & Description \\
    \midrule
    add  & Add value to table value & ln   & Natural log of table value \\
    sub  & Subtract value from table value & log  & Log base 10 of table value \\
    mult & Multiply value by table value & e    & $e$ raised to the table-value power \\
    div  & Divide table value by value & sin  & Sine of table value (in degrees) \\
    idiv & Divide value by table value & cos  & Cosine of table value (in degrees) \\
    exp  & Raise table value to power of value & tan  & Tangent of table value (in degrees) \\
    iexp & Raise value to power of table value & asin & $\sin^{-1}$ of table value (in degrees) \\
    max  & Limit table value to less than a value & acos & $\cos^{-1}$ of table value (in degrees) \\
    min  & Limit table value to greater than a value & atan & $\tan^{-1}$ of table value (in degrees) \\
    set  & Set table value to a value & zero & Set table value to zero \\
    abs  & Absolute value of table value & csto & Store table value and set it to zero \\
    neg  & Negate table value & if   & If-then statement \\
    sqr  & Square table value & goto & Goto statement \\
    sqrt & Square root of table value & end  & End of table \\
    \bottomrule
  \end{tabularx}
\end{table}

Operations that act only on the current table value require no additional constant or
table. For example,

\begin{lstlisting}[style=taosmanual]
start
  add alpha
  sin
  sqr
  add 1.0
end
\end{lstlisting}

computes $1+\sin^2\alpha$. The \texttt{sin} and \texttt{sqr} operations require no
additional values; they act only on the current table value. The \texttt{add}
operation does require another value. Angle of attack is added in the first
\texttt{add} statement, and the constant 1.0 is added in the second.

In \cref{tab:math-operations}, operations \texttt{add} through \texttt{set} require
an additional value; operations \texttt{abs} through \texttt{zero} do not. The
\texttt{csto}, \texttt{if}, and \texttt{goto} operations are special operations
discussed later. An additional value can be a constant, a state variable, a value
interpolated from a table, a temporary storage variable, or a user-defined variable.

\subsection{Constant Values and State Variable Values}\label{sec:table-constant-state-values}

Constant values are always treated as real numbers. They can be entered with or
without decimal points, and scientific notation in e-format is acceptable. For
example,

\begin{lstlisting}[style=taosmanual]
add 2.0
mult 1.5e-6
exp 3
\end{lstlisting}

computes $\bigl((x+2.0)(1.5\times10^{-6})\bigr)^3$, where $x$ is the current table
value.

State-variable values change as the flight path is integrated. When table values are
computed during path integration, current state-variable values can be retrieved by
using the names in \cref{tab:taos-state-variables}. For example, dynamic pressure,
defined as $0.7pM^2$, can be computed with

\begin{lstlisting}[style=taosmanual]
add 0.7
mult pres
mult mach
mult mach
\end{lstlisting}

Default units for these variables are given in the appendix and can be changed with
the \problemblock{*units/fmt} data block (\cref{sec:block-units-fmt}).

\subsection{Tabulated Data Values}\label{sec:tabulated-data-values}

Values used by math operations can be linearly interpolated from a table that is a
function of as many as five independent variables. The tables use the same form as
simple tables. The dependent-variable name is given first, followed by a list of
independent variables in parentheses. Values for each independent variable define a
matrix, and dependent-variable values are then given for every combination of
independent values.

The simplest table is a function of one variable. For example,

\begin{lstlisting}[style=taosmanual]
add cdh(alt)
  alt = 0, 50000, 100000, 150000, 200000
  cdh = 0.0, 0.002, 0.006, 0.011, 0.042
\end{lstlisting}

identifies a value called \texttt{cdh} as a function of altitude. The interpolated value
of \texttt{cdh} is added to the current table value. The name \texttt{cdh} must be
unique: it cannot be a state variable or user-defined variable. It is only an
identification name and is not available as a storage variable for later operations.

In the next example, the current table value is multiplied by a factor obtained from
a tabulated function of Mach number and a user-defined variable called
\texttt{type}:

\begin{lstlisting}[style=taosmanual]
mult factor(type,mach)
  type = 0, 1
  mach = 0.0, 0.999, 1.000, 25.0
  factor = 0.5, 0.5, 1.0, 1.0, # Values for type=0
           1.0, 1.0, 2.0, 2.0  # Values for type=1
\end{lstlisting}

The value of \texttt{type} is set in the problem file, for example in an
\problemblock{*aero} or \problemblock{*prop} data block. If \texttt{type} is zero,
the current table value is multiplied by 0.5 at subsonic Mach numbers and 1.0 at
supersonic Mach numbers. If \texttt{type} is one, those factors are doubled.

Tabulated functions of more than two variables are input using the pattern described
for simple tables in \cref{sec:simple-tables}. A more general skewed format is
described in \cref{sec:skewed-tabulated-data}.

\subsection{Storage Variables}\label{sec:storage-variables}

The \texttt{csto} (clear and store) operation saves a temporary calculation for later
use, like a storage register on a calculator. After a value is saved and named, it can
be retrieved by subsequent operations. For example,

\begin{lstlisting}[style=taosmanual]
add alpha
cos
sqr
csto ca2
\end{lstlisting}

calculates $\cos^2\alpha$ and saves it as the variable \texttt{ca2}. The
\texttt{csto} operation also sets the current table value to zero in preparation for
subsequent operations.

The storage-variable name supplied to \texttt{csto} must be unique; it cannot be a
state variable from \cref{tab:taos-state-variables} or a user-defined variable used
elsewhere in the table or problem files.

The variable \texttt{ca2}, containing $\cos^2\alpha$, can be used in subsequent
operations either as an operation value or as an independent variable of a table. For
example,

\begin{lstlisting}[style=taosmanual]
add alpha
cos
sqr
csto ca2
add factor(ca2,mach)
\end{lstlisting}

uses \texttt{ca2} as an independent variable. This makes it possible to enter tables
that are complex functions of state variables; here, \texttt{factor} is a function of
$\cos^2\alpha$ and Mach number.

\subsection{User-Defined Variables}\label{sec:table-user-defined-variables}

If a variable name in a math operation is neither a state variable nor a storage
variable, TAOS assumes that it is a user-defined variable. Its value must be supplied
in the problem file in an \problemblock{*aero}, \problemblock{*constants},
\problemblock{*cg}, \problemblock{*prop}, or \problemblock{*wind} data block
(\cref{sec:block-aero,sec:block-constants,sec:block-cg,sec:block-prop,sec:block-wind}).
These are the same places where table identification names are given.

User-defined variables can be used in math operations like state variables. They are
constants whose values are set in the problem file. The advantage of using a
user-defined variable instead of a literal table-file constant is that the value can be
automatically surveyed in the problem file (\cref{sec:block-survey}). If a table
constant can take on different values, making it user-defined is usually preferable.
For example,

\begin{lstlisting}[style=taosmanual]
(constant-thr)
  table thrust units=lb
  start
    sub pres
    mult anoz
    add tvac
  end
\end{lstlisting}

uses two user-defined variables, \texttt{anoz} and \texttt{tvac}. Because this is a
thrust table, their values are supplied in the \problemblock{*prop} data block. The
table is generic in the sense that it can represent any constant-thrust rocket motor.

User-defined variables are powerful but must be documented carefully. A table file
with many user-defined variables can be difficult to use later if the variables are not
recorded. Comments should document every required user-defined variable.

\subsection{If-Then Statements}\label{sec:if-then-statements}

The \texttt{if} operation provides a simple way to control the flow of calculations in
a table. An if statement has the form

\begin{lstlisting}[style=taosmanual]
if (relationship) then math-operation
\end{lstlisting}

where a relationship has the form

\begin{lstlisting}[style=taosmanual]
value relation value
\end{lstlisting}

and each value is a constant or variable name. The relation is \texttt{<},
\texttt{>}, or \texttt{=}. For example,

\begin{lstlisting}[style=taosmanual]
if (mach > 5.0) then add f2(time) no-extrap
  time = 0, 1, 2, 3
  f2 = 2.5, 2.6, 2.95, 3.8
\end{lstlisting}

adds the value \texttt{f2}, interpolated as a function of time, if Mach number is
greater than 5.

The relationship is evaluated first. If it is true, the operation following
\texttt{then} is executed. If it is false---in the example, if Mach number is less than
or equal to 5---nothing is done, and processing continues with the next operation.

Only simple relationships are allowed. Combinations such as
\texttt{mach < 5 \& alt > 50000} are not allowed. Separate relations for not equal
to and greater than or equal to are not provided because they are redundant.

\subsection{Goto Statements}\label{sec:goto-statements}

\texttt{goto} operations and labels also control calculation flow and are often used
with \texttt{if} statements. Every math operation can have an optional label. A label
is any name or value placed before an operation and separated from it by a colon:

\begin{lstlisting}[style=taosmanual]
label-1: add time
label-2: sub tmark
lable-3: csto delta-time
\end{lstlisting}

Labels, like other names, cannot contain delimiters such as blanks or commas, so
underscores or dashes can be used to improve readability.

Labels are referenced by \texttt{goto} operations. For example,

\begin{lstlisting}[style=taosmanual]
start
super: if (mach < 1.0) then goto sub # "if" test
       add mach
       mult mach                         # This set of math
       csto mach2                        # operations is for
       add cdsup(mach2)                  # Mach > 1 (supersonic)
         mach2 = 1.0, 4.0, 9.0, 25.0
         cdsup = 0.120, 0.094, 0.075, 0.055
       goto done

sub:   add cdsub(alpha)                  # This set of math
         alpha = 0, 5, 10                # operations is for
         cdsub = 0.020, 0.023, 0.027     # Mach < 1 (subsonic)

done:  end
\end{lstlisting}

performs one set of calculations when Mach number is subsonic and a different set
when it is supersonic.

\subsection{Skewed Tabulated Data}\label{sec:skewed-tabulated-data}

The tabulated sets discussed so far have a dependent value for every combination of
independent values. The independent values form a matrix, and dependent values are
given for every combination. This is shown on the left of
\cref{fig:square-skewed-data} for a function of two independent variables. The
matrix forms a square for two independent variables or a cube for three independent
variables, so these are called \emph{square tables}.

\begin{figure}[htbp]
  \centering
  \begin{tikzpicture}[x=0.72cm,y=0.72cm,font=\scriptsize]
  \newcommand{\basegrid}[1]{%
    \begin{scope}[shift={#1}]
      \draw[line width=0.75pt] (0,0) rectangle (5,5);
      \foreach \x in {1,2,3,4} {\draw[densely dotted] (\x,0)--(\x,5);}
      \foreach \y in {1,2,3,4} {\draw[densely dotted] (0,\y)--(5,\y);}
      \node[below,font=\small] at (2.5,0) {$x$};
      \node[left,font=\small] at (0,2.5) {$f(x,z)$};
      \node[left] at (0.10,3.60) {$z_1$};
      \node[left] at (0.10,2.48) {$z_2$};
      \node[left] at (0.10,1.28) {$z_3$};
    \end{scope}%
  }
  \basegrid{(0,0)}
  \basegrid{(8.0,0)}
  \node[font=\small] at (2.5,5.55) {Square Table};
  \node[font=\small] at (10.5,5.55) {Skewed Table};

  \begin{scope}
    \draw[line width=0.85pt] (0.5,3.45)--(1.5,4.0)--(2.5,4.25)--(3.5,4.2)--(4.5,3.55);
    \draw[line width=0.85pt] (0.5,2.30)--(1.5,2.85)--(2.5,3.10)--(3.5,2.85)--(4.5,2.25);
    \draw[line width=0.85pt] (0.5,1.15)--(1.5,1.65)--(2.5,1.85)--(3.5,1.65)--(4.5,1.30);
    \foreach \x/\a/\b/\c in {0.5/3.45/2.30/1.15,1.5/4.0/2.85/1.65,2.5/4.25/3.10/1.85,3.5/4.2/2.85/1.65,4.5/3.55/2.25/1.30}{
      \draw[dashed] (\x,0.25)--(\x,4.55);
      \draw[fill=white] (\x-0.1,\a-0.1) rectangle ++(0.2,0.2);
      \draw[fill=white] (\x-0.1,\b-0.1) rectangle ++(0.2,0.2);
      \draw[fill=white] (\x-0.1,\c-0.1) rectangle ++(0.2,0.2);
    }
  \end{scope}

  \begin{scope}[shift={(8.0,0)}]
    \draw[line width=0.85pt] (1.2,3.55)--(2.4,4.00)--(3.7,3.85);
    \draw[line width=0.85pt] (0.85,2.40)--(1.8,2.80)--(2.55,3.00)--(3.45,2.75)--(4.35,2.20);
    \draw[line width=0.85pt] (0.55,1.10)--(1.5,1.55)--(2.55,1.75)--(3.8,1.50)--(4.65,0.90);
    \foreach \x/\y in {1.2/3.55,2.4/4.00,3.7/3.85,0.85/2.40,1.8/2.80,2.55/3.00,3.45/2.75,4.35/2.20,0.55/1.10,1.5/1.55,2.55/1.75,3.8/1.50,4.65/0.90}
      {\draw[fill=white] (\x-0.1,\y-0.1) rectangle ++(0.2,0.2);}
    \foreach \x in {1.2,2.4,3.7} {\draw[dashed] (\x,3.2)--(\x,4.35);}
    \foreach \x in {0.85,1.8,2.55,3.45,4.35} {\draw[dashed] (\x,2.0)--(\x,3.25);}
    \foreach \x in {0.55,1.5,2.55,3.8,4.65} {\draw[dashed] (\x,0.55)--(\x,1.95);}
  \end{scope}
\end{tikzpicture}

  \caption{Square and skewed tabulated data.}
  \label{fig:square-skewed-data}
\end{figure}

Sets of tabulated data are often incomplete. The right side of
\cref{fig:square-skewed-data} shows a function defined at different $x$ values for
$z_1$, $z_2$, and $z_3$. Both the values and the number of points differ among the
curves. TAOS calls such data sets \emph{skewed}. They can only be input using the
full-table format; skewed tables are not allowed in the simple-table format.

The following example shows how a skewed set is input:

\begin{lstlisting}[style=taosmanual]
add cxo(alt,mach)
  alt = 0, 50000
  mach = 3, 5, 7, 9
  cxo = .020, .021, .019, .018, # Values for alt=0
        .021, .022, .021, .020  # Values for alt=50000

  alt = 100000, 150000
  mach = 7, 10, 15
  cxo = .024, .023, .021,       # Values for alt=100000
        .027, .025, .024        # Values for alt=150000
\end{lstlisting}

The data is broken into smaller sets that are square in nature. Values at the same
Mach numbers are available for altitudes of 0 and 50,000 feet, so that part is input
first. Values at a different set of Mach numbers are available at altitudes of 100,000
and 150,000 feet, so those are input next. This is similar to entering separate tables
for low and high altitudes.

The same table can be broken down further:

\begin{lstlisting}[style=taosmanual]
add cxo(alt,mach)
  alt = 0
  mach = 3, 5, 7, 9
  cxo = .020, .021, .019, .018

  alt = 50000
  mach = 3, 5, 7, 9
  cxo = .021, .022, .021, .020

  alt = 100000
  mach = 7, 10, 15
  cxo = .024, .023, .021

  alt = 150000
  mach = 7, 10, 15
  cxo = .027, .025, .024
\end{lstlisting}

Here a one-dimensional table is given for each altitude. The one-dimensional tables
can have different Mach-number values and different numbers of points.

In the general case, one-dimensional tables are given for the last, or rightmost,
independent variable for every combination of the remaining independent variables
and values. The sets are grouped or nested so that values for the leftmost
independent variable form the outermost loop. The following three-dimensional
example illustrates the pattern:

\begin{lstlisting}[style=taosmanual]
add ca(alt,mach,alpha)

  alt=10000, mach=5, alpha=0,2,4,6,8 # Values for 10,000
  ca = .20, .21, .23, .26, .28       # start here

  alt=10000, mach=10, alpha=0,2,4,6
  ca = .19, .19, .20, .22

  alt=10000, mach=15, alpha=0,2,4
  ca = .18, .19, .21

  alt=30000, mach=10, alpha=0,2,4,6  # Values for 30,000
  ca = .21, .23, .25, .29            # start here

  alt=30000, mach=15, alpha=0,2,4
  ca = .20, .21, .23
\end{lstlisting}

A one-dimensional angle-of-attack table is given for each combination of Mach
number and altitude, but the number and values of angles of attack can differ for
each table. All independent values for one set are placed on the same line; because
input is free-field, this is allowed.

Altitude is the leftmost independent variable, so altitude values form the outermost
loop and Mach-number values are nested within them. The altitude value must be
repeated until all Mach-number values have been supplied. This example contains
more Mach numbers at 10,000 feet than at 30,000 feet. Altitude, Mach number, and
angle of attack are all independent variables, so their values must each be strictly
increasing or strictly decreasing.
